(1) 选绕环轴同心的圆形安培回路，理想密绕条件下仅在环芯 $a<r<b$、指定高度范围内有 $\varphi$ 向磁场；包围 $N$ 匝电流给 $B_\varphi(r)2\pi r=\mu_0NI$。所以 $B_\varphi(r)=\mu_0NI/(2\pi r)$；在理想环芯之外场取零，实际有限绕组有小漏磁。
(2) 单匝横截面磁通 $\Phi=\int_0^h\!\int_a^bB_\varphi(r)\,dr\,dz=\mu_0NIh\ln(b/a)/(2\pi)$。全部 $N$ 匝磁链为 $N\Phi=\mu_0N^2Ih\ln(b/a)/(2\pi)$，不能将单匝磁通误当整绕组磁链。
(3) 按 $N\Phi=LI$ 得 $L=\mu_0N^2h\ln(b/a)/(2\pi)$。独立积分 $U_B=\int_a^b[B_\varphi^2/(2\mu_0)](2\pi rh)\,dr=\mu_0N^2I^2h\ln(b/a)/(4\pi)=LI^2/2$，两种求法严格一致。
(4) 当 $d=b-a\ll a$，$\ln(b/a)=\ln(1+d/a)\simeq d/a$，故 $L\simeq\mu_0N^2hd/(2\pi a)=\mu_0N^2A_{\rm cross}/(2\pi a)$，其中 $A_{\rm cross}=hd$。$\mu_0$ 的单位为 $\mathrm{H\,m^{-1}}$，再乘以长度得到亨利；当环芯横截面缩小至零时电感也趋零。
